\section{Background}


\subsection{Production Methods}


\subsection{Fuel forms}
\subsection{HALEU in reactors}
A lot of work has been done to evaluate the performance of reactors 
using \gls{HALEU} fuel. Burns et al. investigated the reactor and fuel cycle 
performance of an \gls{LWR} fueled by an enrichment above 5\% \cite{burns_reactor_2020}.
Their work showed that increasing the fuel enrichment up to 7\% does not 
largely affect the fuel temperature or the moderator temperature coefficients,
but it does decrease the soluble boron coefficient and increase the maximum 
burnup at the edges of the fuel pellets. The impacts on the fuel cycle include 
an increase in the amount of natural uranium required, a decrease in the 
amount of high-level waste disposed of per unit energy, and the 
radioactivity of the \gls{SNF} and \gls{HLW} reduces slightly, 100 years 
after discharge \cite{burns_reactor_2020}.

Increasing the fuel enrichment of the NuScale \gls{SMR} has also been 
investigated \cite{carlson_implications_2022}. Increasing the fuel 
enrichment to 8.34\% (compared with the current <5\% design) doubled 
the fuel discharge burnup and the cycle time. 
Fission product poisons, specifically $^{149}$Sm and $^{135}$Xe, increased 
in concentration when using \gls{HALEU} fuel compared with using the base 
design. This increase impacts the 
neutron poisons needed during operation and the radioactivity of 
the fuel after discharge. Using \gls{HALEU} in the core also led to a 
reduction in the \gls{LCOE} of the reactor for certain combinations of 
enrichment and cycle length \cite{carlson_economic_2020,carlson_implications_2022}.

The fabrication of the \gls{HALEU} fuel can also affect the fuel's 
performance in the reactor. Analysis of fuel for the \gls{MIT} 
Reactor showed that manufactured fuel had a slightly higher $^{235}$U 
concentration (19.8\% instead of 19.75\%) and a slightly higher 
density than original models accounted for \cite{mascolino_impact_2022}. 
These changes in the fuel caused a greater $^{235}$U loading in the core, 
increasing the \keff by 431 pcm \cite{mascolino_impact_2022}. The 
differences in modeled and manufactured fuel meant that the modelers 
had to account for more statistical uncertainty in the reactor performance, 
which can be computationally intensive. To alleviate potential 
computational costs, the modelers developed "high reactivity" and 
"low reactivity" fuel compositions based on uncertainty in the manufactured 
fuel compositions. This additional analysis of the extra fuel compositions 
was sufficient to show that the reactor could still meet safety margins 
with the fuel composition variations \cite{mascolino_impact_2022}.

If \gls{HALEU} fuel is produced by downblending \gls{HEU}, it is possible 
the \gls{HALEU} will contain impurities of minor uranium 
isotopes \cite{nuclear_energy_institute_establishing_2022}. Bounding studies on the 
uranium 
isotopic composition of \gls{HALEU} produced from spent fuel from \gls{EBR} 
shows the presence of $^{232}$U, $^{233}$U, $^{234}$U, and $^{236}$U
\cite{vaden_isotopic_2018}. 
\gls{HALEU} created from downbleneded \gls{HEU} at the Y-12 Nuclear Security 
Complex also has $^{232}$U, $^{234}$U, and $^{236}$U impurities
\cite{nelson_foreign_2010}. $^{233}$U is a fissile 
isotope and is expected to affect the neutron population inside a reactor.
$^{232}$U, $^{234}$U, and $^{236}$U are parasitic neutron absorbers, so 
their presence is also expected to affect the neutron population. The 
reactor neutronics performance of using \gls{HALEU} containing $^{235}$U 
and $^{238}$U was compared 
with the performance using \gls{HALEU} from Y-12 with the known impurities 
\cite{celikten_effects_2021}. By performing burnup dependent \gls{MCNP} 
simulations of the potential Replacement Reactor Concept at \gls{NIST}, 
the authors showed that using the \gls{HALEU} with the \gls{HEU} 
impurities showed a 0.43\% decrease in reactivity, leading to an 8\% reduction 
in the cycle length. This analysis suggests that these impurities may 
impact the performance of commercially deployed reactors as well, which 
warrants further investigation of the reactor performance when using 
downblended \gls{HEU} as fuel. 